\subsection{Теория и интерфейсы API}
\subsection{RGB-модель и диапазоны}
\hyperlink{term:rgb}{RGB} — три управляющих числа: \verb|R| (красный), \verb|G| (зелёный), \verb|B| (синий) в \verb|[0.0, 1.0]|. \verb|0.0| — канал выключен, \verb|1.0| — горит ярко. Смешивая каналы, получаем цвета: \verb|{1,0,0}| (красный), \verb|{0,1,0}| (зелёный), \verb|{0,0,1}| (синий), \verb|{1,1,0}| (жёлтый), \verb|{1,1,1}| (белый), \verb|{0,0,0}| (выключено). Плавная смена значений позволяет создавать анимации.
